\chapter{Limits at Infinity and Asymptotes}

So far our limits have asked a local question: as $x$ creeps toward some finite
landmark $a$, where does $f(x)$ head? In this chapter we step back and ask a
global one. \emph{Where does the function go in the long run?} Walk a mile to
the right along the graph, then ten miles, then a thousand. Does the curve
settle down onto a level highway? Does it shoot off to the heavens? Does it
hug a tilted line forever without ever touching it?

This is the study of \textbf{end behavior}, and its natural language is the
\textbf{limit at infinity}. It is also the study of \textbf{asymptotes} ---
the straight lines a curve approaches but (often) never reaches. Asymptotes are
the skeleton of a graph. If you know the asymptotes, you already know the
silhouette of the function before you plot a single careful point. And once you
can read end behavior fluently, you can compare how fast different functions
grow --- which is the single most useful skill in all of applied mathematics,
from sorting algorithms to the stability of bridges to the half-life of a drug.

\chapterline

\section{What ``Limit at Infinity'' Really Means}
\label{ch09:meaning}

When we wrote $\lim_{x\to a} f(x) = L$ for a finite $a$, the symbol $x\to a$
meant ``$x$ gets arbitrarily close to $a$.'' But $\infty$ is not a number, and
$x$ cannot get \emph{close} to it --- there is nothing to get close to. So we
have to say precisely what the new symbol asserts.

\begin{happeningbox}
The statement $\lim_{x\to\infty} f(x) = L$ does \textbf{not} say that $x$
arrives anywhere. It is a promise about the long run:

\begin{center}
\emph{No matter how tight a tolerance you demand around $L$, I can name a
starting point far enough to the right so that, from there on, $f(x)$ stays
inside that tolerance.}
\end{center}

There is no last value of $x$, no destination. There is only the guarantee that
\emph{eventually and forever after}, the output is pinned near $L$. The whole
content of ``$x\to\infty$'' is the word \emph{eventually}.
\end{happeningbox}

Here is the precise version, the exact analogue of the $\eps$--$\delta$
definition but with the role of ``close to $a$'' played by ``far to the right.''

\begin{definition}[Limit at infinity]
\label{ch09:def-limit-inf}
We write $\displaystyle\lim_{x\to\infty} f(x) = L$ if: for every $\eps > 0$
there exists a real number $M$ such that
\[
   x > M \quad\Longrightarrow\quad \abs{f(x) - L} < \eps .
\]
Likewise $\displaystyle\lim_{x\to-\infty} f(x) = L$ if for every $\eps>0$ there
is an $M$ with $x < M \implies \abs{f(x)-L}<\eps$.
\end{definition}

Compare this with the finite-$a$ definition. There, $\eps$ bought you a $\delta$
and you looked inside the window $0<\abs{x-a}<\delta$. Here, $\eps$ buys you a
threshold $M$ and you look at the entire tail $x>M$. The picture is not a
shrinking window around a point; it is a horizontal corridor of half-width
$\eps$ around the height $L$, and the claim is that the graph enters that
corridor and never leaves.

\begin{figure}[H]
\centering
\begin{tikzpicture}
\begin{axis}[
  width=11cm, height=6cm,
  xlabel=$x$, ylabel=$f(x)$,
  domain=0.5:10, samples=120,
  ymin=-0.2, ymax=3.2, xmin=0, xmax=10,
  axis lines=left,
  xtick={6}, xticklabels={$M$},
  ytick={2}, yticklabels={$L$},
  clip=false,
]
\addplot[draw=none, name path=top] {2.25};
\addplot[draw=none, name path=bot] {1.75};
\addplot[areafill!40] fill between[of=top and bot, soft clip={domain=6:10}];
\addplot[curveblue, very thick] {2 + 1.4/(x)};
\draw[dashed, gray] (axis cs:0,2) -- (axis cs:10,2);
\draw[dashed, trapred] (axis cs:6,-0.2) -- (axis cs:6,3.2);
\node[anchor=west, trapred] at (axis cs:6.1,3.0) {tail $x>M$};
\node[anchor=west] at (axis cs:6.2,2.45) {corridor $L\pm\eps$};
\end{axis}
\end{tikzpicture}
\caption{A limit at infinity. Once $x$ passes the threshold $M$, the graph is
trapped inside the corridor of half-width $\eps$ around the height $L$. Shrink
$\eps$ and you may need a larger $M$, but such an $M$ always exists.}
\label{ch09:fig-corridor}
\end{figure}

\begin{labbox}
Show directly from Definition~\ref{ch09:def-limit-inf} that
$\lim_{x\to\infty} \tfrac{1}{x} = 0$. Given $\eps>0$, we need $x>M \implies
\abs{1/x - 0} < \eps$, i.e. $1/x < \eps$ (for $x>0$), i.e. $x > 1/\eps$. So
\emph{choose $M = 1/\eps$.} Then any $x>M$ satisfies $1/x < 1/M = \eps$. Done.
Notice how $M$ depends on $\eps$: demand a tolerance ten times tighter and your
threshold moves ten times farther out. That is the whole mechanism.
\end{labbox}

\subsection{Infinite limits: when the output runs away}

There is a second, completely different way for a function to behave badly: the
\emph{output} can blow up. We already met this near vertical features like
$1/x$ at $x=0$. Let us pin down the language, because ``$=\infty$'' is another
phrase that needs a careful reading.

\begin{definition}[Infinite limit]
\label{ch09:def-inf-limit}
We write $\displaystyle\lim_{x\to a} f(x) = +\infty$ if for every $N>0$ there is
a $\delta>0$ such that
\[
  0 < \abs{x-a} < \delta \quad\Longrightarrow\quad f(x) > N .
\]
That is: name any height $N$, however large, and the graph eventually climbs
above it and stays above it as $x$ closes in on $a$. The case $-\infty$ and the
cases $x\to\pm\infty$ are defined by the obvious changes.
\end{definition}

\begin{trapbox}
The equals sign in $\lim f(x) = \infty$ is a \textbf{convenient lie}. There is
no number $\infty$ that the limit equals. Strictly, the limit \emph{does not
exist} --- we are just describing \emph{how} it fails to exist: not by
oscillating, not by jumping, but by growing without bound. Treat ``$=\infty$''
as one indivisible phrase meaning ``increases beyond every bound,'' never as
an arithmetic value you can add or divide.
\end{trapbox}

\begin{historybox}
The struggle to make these ideas rigorous is the great drama of $19$th-century
analysis. Newton and Leibniz computed with quantities that were ``ultimately''
vanishing or ``infinitely'' large and got correct answers by sure instinct, but
Bishop Berkeley mocked the infinitesimals as ``ghosts of departed quantities,''
and he had a point: nobody could say what they \emph{were}. The repair came in
stages. Bernard Bolzano (1810s) and especially Augustin-Louis Cauchy (1820s)
recast limits as statements about getting arbitrarily close. Then Karl
Weierstrass, lecturing in Berlin in the 1860s, supplied the final grammar --- the
$\eps$, the $\delta$, and (for our chapter) the threshold $M$ --- so that
``tends to infinity'' became a finite, checkable sentence with no ghosts in it.
Definition~\ref{ch09:def-limit-inf} is Weierstrass's gift: a way to talk about
the infinite using only finite quantities.
\end{historybox}

\section{Horizontal Asymptotes: the Level Highway}
\label{ch09:horizontal}

A horizontal asymptote is simply a finite limit at infinity, read off the graph.

\begin{definition}[Horizontal asymptote]
\label{ch09:def-ha}
The line $y=L$ is a \textbf{horizontal asymptote} of $f$ if
$\lim_{x\to\infty} f(x) = L$ or $\lim_{x\to-\infty} f(x) = L$ (or both).
\end{definition}

A few things worth saying out loud, because each contradicts a common belief.

\begin{itemize}
\item A graph \textbf{may cross its horizontal asymptote}, even infinitely
often. The asymptote describes the \emph{tail}, not a barrier. The function
$f(x) = \dfrac{\sin x}{x}$ crosses $y=0$ at every nonzero multiple of $\pi$ yet still
has $y=0$ as a horizontal asymptote, because the crossings get ever smaller.
\item A function can have \textbf{two different} horizontal asymptotes, one to
the right and one to the left. We will meet $\arctan x$ and the logistic curve,
which level off at different heights in the two directions.
\item Having a horizontal asymptote is a statement about ends only. It says
nothing about what happens in the middle, where the function may do anything.
\end{itemize}

\begin{figure}[H]
\centering
\begin{tikzpicture}
\begin{axis}[
  width=12cm, height=5.2cm,
  xlabel=$x$, domain=-12:12, samples=200,
  ymin=-1.4, ymax=1.4, xmin=-12, xmax=12,
  axis lines=middle, ytick={-1,1}, xtick=\empty,
  clip=false,
]
\addplot[curveblue, very thick] {sin(deg(x))/x};
\draw[dashed, trapred] (axis cs:-12,0) -- (axis cs:12,0);
\node[anchor=south west, trapred] at (axis cs:7,0.02) {$y=0$};
\end{axis}
\end{tikzpicture}
\caption{$f(x)=\sin(x)/x$ crosses its horizontal asymptote $y=0$ at every nonzero
multiple of $\pi$. The asymptote governs the \emph{tail}, not a forbidden line.
The wiggles decay because the amplitude $1/\abs{x}$ shrinks.}
\label{ch09:fig-sinc}
\end{figure}

\subsection{Two asymptotes, one curve: the inverse tangent and the logistic}

Some of the most important curves in science have a floor on one side and a
ceiling on the other. The inverse tangent levels off at $\pi/2$ as $x\to+\infty$
and at $-\pi/2$ as $x\to-\infty$. The logistic function
$\sigma(x) = \dfrac{1}{1+e^{-x}}$ --- the workhorse of population models and of
machine-learning classifiers --- rises from the floor $y=0$ to the ceiling
$y=1$. In both cases the graph is sandwiched between two horizontal highways.

\begin{figure}[H]
\centering
\begin{tikzpicture}
\begin{axis}[
  width=12cm, height=5.4cm,
  xlabel=$x$, domain=-8:8, samples=160,
  ymin=-0.15, ymax=1.15, xmin=-8, xmax=8,
  axis lines=middle, ytick={0,0.5,1}, xtick=\empty,
  clip=false, legend pos=south east, legend style={font=\small},
]
\addplot[curveblue, very thick] {1/(1+exp(-x))};
\addlegendentry{$\sigma(x)=\frac{1}{1+e^{-x}}$}
\draw[dashed, trapred] (axis cs:-8,1) -- (axis cs:8,1);
\draw[dashed, trapred] (axis cs:-8,0) -- (axis cs:8,0);
\node[anchor=west, trapred] at (axis cs:5,0.95) {$y=1$ ceiling};
\node[anchor=west, trapred] at (axis cs:-7.8,0.06) {$y=0$ floor};
\end{axis}
\end{tikzpicture}
\caption{The logistic curve $\sigma(x)$ has two horizontal asymptotes: a floor
at $y=0$ (as $x\to-\infty$) and a ceiling at $y=1$ (as $x\to+\infty$). Most
saturating processes in nature look like this.}
\label{ch09:fig-logistic}
\end{figure}

\begin{applicationbox}
\textbf{Why classifiers love the ceiling and floor.} A neural network outputs a
raw score $z$ that can be any real number, but a probability must live in
$[0,1]$. Passing $z$ through $\sigma$ squashes it: huge positive $z$ saturates
near $1$ (``confident yes''), huge negative $z$ saturates near $0$ (``confident
no''). The two horizontal asymptotes are precisely what make $\sigma$ a sensible
probability. The same shape models a population approaching its carrying
capacity, a chemical reaction approaching completion, or a rumor saturating a
town --- anything that grows fast at first and then runs out of room.
\end{applicationbox}

\section{Vertical Asymptotes: the Cliff}
\label{ch09:vertical}

A horizontal asymptote happens because $x$ ran off to infinity. A
\textbf{vertical} asymptote is the opposite: $x$ stays at a finite spot, but the
\emph{output} runs to infinity.

\begin{definition}[Vertical asymptote]
\label{ch09:def-va}
The line $x=a$ is a \textbf{vertical asymptote} of $f$ if at least one of the
one-sided limits $\lim_{x\to a^+} f(x)$ or $\lim_{x\to a^-} f(x)$ is $+\infty$
or $-\infty$.
\end{definition}

The usual culprit is a denominator that hits zero while the numerator does not.
Near such a point the fraction has a fixed, nonzero amount on top divided by a
vanishing amount on the bottom, and ``fixed over tiny'' is ``huge.'' The sign on
each side tells you whether the cliff goes up or down.

\begin{labbox}
Find the vertical asymptote of $f(x)=\dfrac{x+2}{x-3}$ and its two-sided
behavior. The denominator vanishes at $x=3$, where the numerator is $5\ne 0$, so
$x=3$ is a vertical asymptote. Now test the signs. Just to the \emph{right} of
$3$, say $x=3.01$: top $\approx 5>0$, bottom $=0.01>0$, ratio $\approx +500$, so
$f\to+\infty$. Just to the \emph{left}, $x=2.99$: top $\approx 5$, bottom
$=-0.01<0$, ratio $\approx -500$, so $f\to-\infty$. The curve dives down on the
left of the cliff and rockets up on the right.
\end{labbox}

\begin{trapbox}
\textbf{A zero in the denominator does not guarantee a vertical asymptote.}
If the numerator vanishes there too, the factor may cancel and leave only a
\emph{hole}. Consider $g(x)=\dfrac{x^2-9}{x-3}=\dfrac{(x-3)(x+3)}{x-3}$. For
$x\ne 3$ this equals $x+3$, a perfectly tame line with a single missing point at
$(3,6)$. No cliff --- just a puncture. Always factor and cancel \emph{before}
declaring an asymptote. Asymptote when the zero survives in the bottom; hole
when it cancels.
\end{trapbox}

\begin{figure}[H]
\centering
\begin{tikzpicture}
\begin{axis}[
  width=12cm, height=6cm,
  xlabel=$x$, ylabel=$y$,
  domain=-4:2.6, samples=120,
  ymin=-12, ymax=12, xmin=-4, xmax=10,
  axis lines=middle, xtick={3}, ytick={1}, xticklabels={$3$},
  yticklabels={$1$}, clip=true, restrict y to domain=-40:40,
]
\addplot[curveblue, very thick, domain=-4:2.6] {(x+2)/(x-3)};
\addplot[curveblue, very thick, domain=3.4:10] {(x+2)/(x-3)};
\draw[dashed, trapred] (axis cs:3,-12) -- (axis cs:3,12);
\draw[dashed, targetgreen] (axis cs:-4,1) -- (axis cs:10,1);
\node[anchor=south east, trapred] at (axis cs:3,9) {$x=3$};
\node[anchor=south west, targetgreen] at (axis cs:6,1.2) {$y=1$};
\end{axis}
\end{tikzpicture}
\caption{$f(x)=\frac{x+2}{x-3}$ has a vertical asymptote at $x=3$ (output runs
off) and a horizontal asymptote at $y=1$ (input runs off). The two kinds of
asymptote are dual: one freezes $x$ and lets $y$ escape, the other freezes the
level and lets $x$ escape.}
\label{ch09:fig-rational}
\end{figure}

\section{Dominant-Term Reasoning for Rational Functions}
\label{ch09:dominant}

How did we know the horizontal asymptote in Figure~\ref{ch09:fig-rational} was
$y=1$? There is a fast, intuitive method that works for every rational function
and, once understood, frees you from memorizing a table of cases.

\begin{intuitionbox}
When $x$ is enormous, a polynomial is \textbf{ruled by its highest-degree term}.
In $3x^2 + 50x + 9000$, plug in $x=10^6$: the $3x^2$ term is $3\times 10^{12}$,
while $50x = 5\times 10^7$ and the constant $9000$ are rounding error by
comparison. The leading term is the elephant; the rest are fleas. So for end
behavior, \emph{keep only the leading term of the top and the leading term of
the bottom, and simplify.}
\end{intuitionbox}

\begin{whybox}
\textbf{Why dominant-term reasoning is exact, not just sloppy.} Take
$f(x)=\dfrac{3x^2+50x+9000}{x^2-7}$. Divide top and bottom by $x^2$, the highest
power present:
\[
  f(x) = \frac{3 + \dfrac{50}{x} + \dfrac{9000}{x^2}}
              {1 - \dfrac{7}{x^2}} .
\]
This is not an approximation; it is an exact equality (for $x\ne 0$). Now every
term with an $x$ in the denominator $\to 0$ as $x\to\infty$, by the
$1/x\to 0$ fact we proved by hand. So the whole thing tends to
$\dfrac{3+0+0}{1-0} = 3$. The ``keep the leading terms'' shortcut is just this
division, performed in your head. Dividing by the dominant power is the
\emph{mechanism}; the shortcut is its shadow.
\end{whybox}

The general rule falls out of comparing the degree of the numerator $p$ with the
degree of the denominator $q$. Let $\deg p = n$ with leading coefficient $a$, and
$\deg q = m$ with leading coefficient $b$.

\begin{focusbox}
\textbf{End behavior of $\dfrac{p(x)}{q(x)}$ as $x\to\pm\infty$.}
\begin{itemize}
\item \textbf{Bottom wins} ($n<m$): the fraction $\to 0$. Horizontal asymptote
$y=0$. (Top grows slower; the big denominator crushes it.)
\item \textbf{Tie} ($n=m$): the fraction $\to a/b$, the ratio of leading
coefficients. Horizontal asymptote $y=a/b$.
\item \textbf{Top wins} ($n>m$): the fraction $\to \pm\infty$. No horizontal
asymptote. If $n=m+1$ exactly, there is instead a \emph{slant} asymptote
(next section).
\end{itemize}
\end{focusbox}

\begin{example}[All three cases at once]
\label{ch09:ex-three-cases}
\[
\frac{2x+1}{x^2+5}\ \to\ 0
\quad(n=1<m=2),
\qquad
\frac{4x^2-x}{2x^2+3}\ \to\ \frac{4}{2}=2
\quad(n=m=2),
\]
\[
\frac{x^3}{x^2+1}\ \to\ \pm\infty
\quad(n=3>m=2).
\]
For the middle one, the horizontal asymptote is $y=2$ in \emph{both} directions,
since the leading coefficients do not care about the sign of $x$ when both
degrees are even and equal.
\end{example}

\begin{trapbox}
When degrees \emph{tie}, students sometimes answer ``$1$'' by reflex, as if the
leading $x$'s simply cancel. They do not: you keep the \emph{coefficients}. In
$\frac{4x^2-x}{2x^2+3}$ the limit is $4/2=2$, not $1$. And the lower-degree
terms ($-x$ and $+3$) are genuinely irrelevant to the limit --- but only to the
\emph{limit}. They still bend the curve before it settles.
\end{trapbox}

\begin{walkbox}
Without computing, decide the end behavior of each, just by eyeing leading
terms: $\dfrac{x^4-100x}{3x^4+1}$, $\dfrac{7x}{x^3-x}$,
$\dfrac{x^2-1}{x+1}$. (Answers: $\to \tfrac13$; $\to 0$; the last one cancels to
$x-1$, so it heads to $\pm\infty$ along a line --- a slant asymptote in disguise.)
\end{walkbox}

\subsection{Roots and other non-polynomial tops}

Dominant-term reasoning is not just for polynomials over polynomials. Whenever
you can identify the fastest-growing piece on top and on bottom, the same divide
trick works --- you just have to be careful about signs under square roots.

\begin{example}[A square root that needs care]
\label{ch09:ex-sqrt}
Find $\lim_{x\to\pm\infty}\dfrac{\sqrt{4x^2+1}}{x+3}$.
The dominant piece on top is $\sqrt{4x^2}=2\abs{x}$, and on the bottom it is $x$.
Divide top and bottom by $x$. The subtlety: $\sqrt{4x^2+1}/x = \pm\sqrt{4+1/x^2}$,
where the sign is $+$ if $x>0$ and $-$ if $x<0$, because $\sqrt{x^2}=\abs{x}$.
Hence
\[
  x\to+\infty:\quad \frac{\sqrt{4+1/x^2}}{1+3/x}\to\frac{\sqrt 4}{1}=2,
  \qquad
  x\to-\infty:\quad \frac{-\sqrt{4+1/x^2}}{1+3/x}\to -2 .
\]
Two different horizontal asymptotes, $y=2$ and $y=-2$, born entirely from the
$\sqrt{x^2}=\abs{x}$ sign rule.
\end{example}

\begin{trapbox}
$\sqrt{x^2}=\abs{x}$, \textbf{not} $x$. For $x\to-\infty$ this single fact flips
a sign and gives you the second asymptote. Forgetting it is the most common
error in limits-at-infinity with radicals. When dividing by $x$ inside a square
root, you are really dividing by $\sqrt{x^2}=\abs{x}$, so a minus sign appears on
the negative side.
\end{trapbox}

\begin{example}[A sneaky difference of roots]
\label{ch09:ex-diff-roots}
Find $\lim_{x\to\infty}\bigl(\sqrt{x^2+x}-x\bigr)$. As written this is
``$\infty-\infty$,'' which decides nothing. Multiply by the conjugate over
itself:
\[
\sqrt{x^2+x}-x
= \frac{(\sqrt{x^2+x}-x)(\sqrt{x^2+x}+x)}{\sqrt{x^2+x}+x}
= \frac{x^2+x-x^2}{\sqrt{x^2+x}+x}
= \frac{x}{\sqrt{x^2+x}+x}.
\]
Now divide top and bottom by $x$ (using $\sqrt{x^2+x}/x=\sqrt{1+1/x}$ for $x>0$):
\[
= \frac{1}{\sqrt{1+1/x}+1}\ \longrightarrow\ \frac{1}{\sqrt 1 + 1}=\frac12 .
\]
The two huge quantities differ by a finite gap of $\tfrac12$ in the limit. The
conjugate trick converts a hopeless $\infty-\infty$ into a clean ratio where
dominant terms can speak.
\end{example}

\section{Slant (Oblique) Asymptotes: the Tilted Highway}
\label{ch09:slant}

When the top of a rational function out-grows the bottom by exactly one degree,
the curve does not flatten and does not merely blow up vertically. Instead it
approaches a \emph{slanted line}. The graph rides a tilted highway off to
infinity.

\begin{intuitionbox}
A slant asymptote is what polynomial \emph{long division} sees. Dividing
$p(x)$ by $q(x)$ when $\deg p = \deg q + 1$ gives a linear quotient $mx+b$ plus a
proper remainder $r(x)/q(x)$:
\[
   \frac{p(x)}{q(x)} = \underbrace{mx+b}_{\text{the line}}
   + \underbrace{\frac{r(x)}{q(x)}}_{\to\, 0}.
\]
The remainder term has a smaller-degree top than bottom, so by the ``bottom
wins'' rule it vanishes at infinity. What is left in the long run is exactly the
line $y=mx+b$. \textbf{The quotient of long division is the asymptote.}
\end{intuitionbox}

\begin{example}[Finding a slant asymptote by division]
\label{ch09:ex-slant}
Consider $f(x)=\dfrac{x^2+2x-1}{x-1}$. Degree $2$ over degree $1$: expect a
slant. Long-divide:
\[
   x^2+2x-1 = (x-1)(x+3) + 2,
   \qquad\text{so}\qquad
   f(x) = x+3 + \frac{2}{x-1}.
\]
As $x\to\pm\infty$ the remainder $\dfrac{2}{x-1}\to 0$, so $f(x)$ approaches the
line $y=x+3$. The sign of $\dfrac{2}{x-1}$ also tells you \emph{which side} of
the line the curve rides: positive (above the line) for $x>1$, negative (below)
for $x<1$. There is also a vertical asymptote at $x=1$, where the remainder
itself blows up.
\end{example}

\begin{figure}[H]
\centering
\begin{tikzpicture}
\begin{axis}[
  width=12cm, height=6.6cm,
  xlabel=$x$, ylabel=$y$,
  domain=-8:0.3, samples=120,
  ymin=-12, ymax=16, xmin=-8, xmax=10,
  axis lines=middle, xtick={1}, ytick=\empty, xticklabels={$1$},
  clip=true, restrict y to domain=-40:40,
]
\addplot[curveblue, very thick, domain=-8:0.6] {(x^2+2*x-1)/(x-1)};
\addplot[curveblue, very thick, domain=1.4:10] {(x^2+2*x-1)/(x-1)};
\addplot[dashed, structurepurple, very thick, domain=-8:10] {x+3};
\draw[dashed, trapred] (axis cs:1,-12) -- (axis cs:1,16);
\node[anchor=south west, structurepurple] at (axis cs:5,7.2) {$y=x+3$};
\node[anchor=south east, trapred] at (axis cs:1,12) {$x=1$};
\end{axis}
\end{tikzpicture}
\caption{$f(x)=\frac{x^2+2x-1}{x-1}$ has a slant asymptote $y=x+3$ and a vertical
asymptote $x=1$. Far out, the curve is indistinguishable from the tilted line;
the gap between them is the remainder $2/(x-1)$, which melts away.}
\label{ch09:fig-slant}
\end{figure}

\begin{whybox}
\textbf{Why ``the quotient is the asymptote'' is a theorem, not a coincidence.}
Two functions $f$ and $g$ are \emph{asymptotic} if $f(x)-g(x)\to 0$ as
$x\to\infty$ --- their graphs become indistinguishable. Long division writes
$f = (mx+b) + \text{remainder}$ with remainder $\to 0$, which is \emph{exactly}
the statement $f(x)-(mx+b)\to 0$. So the line $mx+b$ is asymptotic to $f$ by
definition. The same logic shows the asymptote need not be a line at all: if the
top exceeds the bottom by two degrees, the quotient is a \emph{parabola}, and the
curve hugs that parabola at infinity. People rarely name these ``parabolic
asymptotes,'' but the mechanism is identical --- whatever the polynomial quotient
is, the curve becomes it.
\end{whybox}

\begin{trapbox}
A graph has a slant asymptote \textbf{or} a horizontal asymptote, never both in
the same direction --- a line is either tilted or level. If $\deg p = \deg q$ you
get a horizontal asymptote; if $\deg p = \deg q + 1$ you get a slant one; if the
gap is $2$ or more there is no straight-line asymptote at all. Always compare
degrees \emph{first} to know which question you are even asking.
\end{trapbox}

\section{Growth Rates Revisited: Who Wins the Long Race?}
\label{ch09:growth}

Dominant-term reasoning settled contests \emph{within} the polynomials. But the
deepest and most useful comparisons are \emph{across} families: logarithms
versus powers versus exponentials. These three grow at wildly different speeds,
and the pecking order never changes once $x$ is large enough.

\begin{focusbox}
\textbf{The growth hierarchy (for large $x$).}
\[
   \ln x \ \ll\ x^{p}\ (p>0) \ \ll\ a^{x}\ (a>1) \ \ll\ x! \ \ll\ x^{x}.
\]
Here $f \ll g$ means $\dfrac{f(x)}{g(x)}\to 0$ as $x\to\infty$: $f$ is utterly
outrun by $g$. \emph{Logs crawl, powers walk, exponentials sprint, factorials
fly.} Every exponential, no matter how gentle its base (say $1.000001^x$),
eventually overtakes every power, no matter how steep (say $x^{1000}$).
\end{focusbox}

The claim is important enough, and counterintuitive enough, to earn a proof
rather than a slogan. Here is \emph{why} the exponential always wins.

\begin{intuitionbox}
\textbf{Why exponentials beat powers: constant percentage versus constant
amount.} A power like $x^2$ grows by an ever-smaller \emph{percentage} as $x$
climbs --- going from $100$ to $101$ barely budges $x^2$ in relative terms. An
exponential like $2^x$ grows by a \emph{fixed percentage} ($100\%$ per unit
step) forever. Fixed-percentage growth compounds; eventually-shrinking-percentage
growth cannot keep up. The race is not close: it is a rout, just delayed.
\end{intuitionbox}

\begin{theorem}[Exponential dominates every power]
\label{ch09:thm-exp-beats-power}
For every fixed $p>0$ and every $a>1$,
\[
   \lim_{x\to\infty} \frac{x^{p}}{a^{x}} = 0 .
\]
\end{theorem}

\begin{proof}
The cleanest path is through logarithms, which turn the lopsided ratio into a
manageable difference. Take logs of the ratio:
\[
  \ln\!\left(\frac{x^{p}}{a^{x}}\right)
  = p\ln x - x\ln a .
\]
We will show this $\to -\infty$; then the ratio itself, being $e$ to that power,
$\to 0$. Factor out $x$:
\[
  p\ln x - x\ln a = x\left( \frac{p\ln x}{x} - \ln a\right).
\]
Here $\ln a > 0$ is a fixed positive constant. So everything hinges on the term
$\dfrac{\ln x}{x}$. If we can show $\dfrac{\ln x}{x}\to 0$, then the bracket
tends to $0-\ln a = -\ln a < 0$, while the factor $x\to+\infty$; a positive-going
$x$ times a negative-bounded bracket gives $-\infty$. So it remains only to
prove the cornerstone fact: $\ln x / x \to 0$.

\textbf{Cornerstone: $\ln x \ll x$.} For $t\ge 1$ we have $\dfrac{1}{t}\le
\dfrac{1}{\sqrt t}$, so integrating from $1$ to $x$,
\[
  \ln x = \int_1^x \frac{1}{t}\,\dd t \le \int_1^x \frac{1}{\sqrt t}\,\dd t
  = 2\sqrt x - 2 \le 2\sqrt x .
\]
Therefore $\dfrac{\ln x}{x} \le \dfrac{2\sqrt x}{x} = \dfrac{2}{\sqrt x}\to 0$.
That bound --- $\ln x \le 2\sqrt x$ --- is the whole secret: the logarithm is
trapped beneath a mere square root, so it is hopelessly slow. Feeding this back,
$\dfrac{p\ln x}{x}\to 0$, the bracket $\to -\ln a<0$, the product $\to -\infty$,
and the original ratio $\to 0$.
\end{proof}

\begin{happeningbox}
\textbf{What the logarithm trick is really doing.} Exponentials and powers live
on different scales, and dividing them directly is like comparing a sprinter's
position to a walker's --- both numbers race off to infinity and the ratio is a
$\infty/\infty$ blur. The logarithm \emph{changes the ruler}: it converts
multiplicative growth into additive growth, $a^x$ into the straight line
$x\ln a$, and $x^p$ into the slow curve $p\ln x$. On this new ruler the contest
is between a \emph{line} and a \emph{logarithm}, and the line obviously wins.
Taking logs is the standard move whenever you must compare two runaway
quantities: it tames both at once.
\end{happeningbox}

\begin{connectionbox}
\textbf{Big-O and the speed of algorithms.} Computer science lives on this
hierarchy. An algorithm that does $\ln n$ work (binary search) is gloriously
fast; $n$ or $n\ln n$ (a good sort) is fine; $n^2$ (a naive sort) is sluggish on
big inputs; and $2^n$ (brute-forcing every subset) is hopeless past a few dozen
items. The symbol ``$f = O(g)$'' means $f$ grows \emph{no faster than} a constant
times $g$ --- a bounded ratio --- and ``$f = o(g)$'' is exactly our $f\ll g$.
The same limits that shape a curve's tail decide whether a program finishes
before the sun burns out.
\end{connectionbox}

\begin{labbox}
\textbf{Feel the rout with your own arithmetic.} Compare $x^{10}$ against $2^x$
at a few points. At $x=10$: $x^{10}=10^{10}=10{,}000{,}000{,}000$ while
$2^{10}=1024$ --- the power is ahead by seven orders of magnitude. At $x=60$:
$2^{60}\approx 1.15\times 10^{18}$ while $60^{10}\approx 6.0\times 10^{17}$ --- the
exponential has just pulled even. Past there, $2^x$ leaves $x^{10}$ in the dust
forever. The crossover is late, but it is inevitable. That delay is exactly why
the limit can be ``obviously'' one thing for small inputs and the \emph{opposite}
thing in truth.
\end{labbox}

\begin{example}[Putting the hierarchy to work]
\label{ch09:ex-hierarchy}
Evaluate each by spotting the dominant family.
\[
\lim_{x\to\infty}\frac{x^{100}}{e^x}=0
\quad(\text{exp beats power}),
\qquad
\lim_{x\to\infty}\frac{\ln x}{x^{0.001}}=0
\quad(\text{power beats log}),
\]
\[
\lim_{x\to\infty}\frac{3^x + x^5}{3^x - 100}
= \lim_{x\to\infty}\frac{1 + x^5/3^x}{1 - 100/3^x}
= \frac{1+0}{1-0}=1 .
\]
In the last one we divided by the dominant term $3^x$, then used
$x^5/3^x\to 0$ (exp beats power) and $100/3^x\to 0$. Same dominant-term
strategy as for rational functions --- only now the ``leading term'' is the
fastest-growing \emph{family}, not the highest power.
\end{example}

\begin{figure}[H]
\centering
\begin{tikzpicture}
\begin{axis}[
  width=12cm, height=6.4cm,
  xlabel=$x$, ylabel=$y$,
  domain=0.3:8, samples=160,
  ymin=0, ymax=60, xmin=0, xmax=8,
  axis lines=left, legend pos=north west, legend style={font=\small},
]
\addplot[targetgreen, very thick] {ln(x)};   \addlegendentry{$\ln x$ (crawls)}
\addplot[curveblue, very thick] {x^2};        \addlegendentry{$x^2$ (walks)}
\addplot[tangentred, very thick] {exp(x)};    \addlegendentry{$e^x$ (sprints)}
\end{axis}
\end{tikzpicture}
\caption{The growth hierarchy on one set of axes. For small $x$ the order can
look ambiguous, but $e^x$ soon rockets past $x^2$, which long ago left $\ln x$
behind crawling along the floor. Extend the window and the gaps only widen.}
\label{ch09:fig-hierarchy}
\end{figure}

\section{Reading a Graph From Its Asymptotes}
\label{ch09:reading}

Put the pieces together. Given a rational function (or any function whose tails
you can analyze), a short checklist lets you sketch an accurate silhouette before
plotting careful points. Asymptotes are the scaffolding; the rest is filling in.

\begin{focusbox}
\textbf{Asymptote checklist for $f(x)=p(x)/q(x)$.}
\begin{enumerate}
\item \textbf{Factor} top and bottom. Cancel common factors --- each cancellation
is a \emph{hole}, not an asymptote.
\item \textbf{Vertical asymptotes:} the surviving zeros of the denominator. Check
the sign on each side to see if the cliff goes up or down.
\item \textbf{Horizontal/slant:} compare degrees. $n<m\Rightarrow y=0$;
$n=m\Rightarrow y=a/b$; $n=m+1\Rightarrow$ slant (do the division); $n>m+1
\Rightarrow$ no line asymptote.
\item \textbf{Intercepts:} $x$-intercepts are the surviving zeros of the
\emph{numerator}; the $y$-intercept is $f(0)$ if defined.
\item \textbf{Sketch} the asymptotes as dashed lines first, then thread the curve
through the intercepts and into the correct corners.
\end{enumerate}
\end{focusbox}

\begin{example}[A full reading]
\label{ch09:ex-full}
Sketch $f(x)=\dfrac{x^2-4}{x^2-1}$.
\emph{Factor:} $f(x)=\dfrac{(x-2)(x+2)}{(x-1)(x+1)}$ --- no common factors, so no
holes. \emph{Vertical asymptotes:} $x=1$ and $x=-1$. \emph{Horizontal:} degrees
tie ($n=m=2$) with leading coefficients $1$ and $1$, so $y=1$. \emph{Intercepts:}
zeros of the numerator give $x=\pm 2$; and $f(0)=\dfrac{-4}{-1}=4$, the
$y$-intercept. With two vertical cliffs, a level highway $y=1$ at both ends, and
the curve pinned through $(\pm2,0)$ and $(0,4)$, the entire shape is forced. The
middle branch has a local minimum at $(0,4)$, staying strictly above the asymptote $y=1$.
\end{example}

\begin{figure}[H]
\centering
\begin{tikzpicture}
\begin{axis}[
  width=12cm, height=6.6cm,
  xlabel=$x$, ylabel=$y$,
  domain=-4:4, samples=200,
  ymin=-6, ymax=8, xmin=-4, xmax=4,
  axis lines=middle, xtick={-2,2}, ytick={1,4},
  clip=true, restrict y to domain=-30:30,
]
\addplot[curveblue, very thick, domain=-4:-1.05] {(x^2-4)/(x^2-1)};
\addplot[curveblue, very thick, domain=-0.95:0.95] {(x^2-4)/(x^2-1)};
\addplot[curveblue, very thick, domain=1.05:4] {(x^2-4)/(x^2-1)};
\draw[dashed, trapred] (axis cs:1,-6) -- (axis cs:1,8);
\draw[dashed, trapred] (axis cs:-1,-6) -- (axis cs:-1,8);
\draw[dashed, targetgreen] (axis cs:-4,1) -- (axis cs:4,1);
\node[anchor=west, targetgreen] at (axis cs:2.3,1.4) {$y=1$};
\node[anchor=south, trapred] at (axis cs:1,6) {$x=1$};
\node[anchor=south, trapred] at (axis cs:-1,6) {$x=-1$};
\end{axis}
\end{tikzpicture}
\caption{$f(x)=\frac{x^2-4}{x^2-1}$, read entirely from its asymptotes:
verticals at $x=\pm1$, a horizontal $y=1$ at both ends, $x$-intercepts at
$\pm2$, and the local minimum $(0,4)$. The skeleton determines the silhouette.}
\label{ch09:fig-full}
\end{figure}

\begin{aestheticbox}
Notice the economy. We never solved $f(x)=k$ for a single generic $k$, never
built a table of twenty points. We asked three structural questions --- where
does the denominator die, how do the degrees compare, where are the zeros --- and
the graph fell out almost for free. This is the recurring pleasure of calculus:
the right \emph{questions} replace mountains of arithmetic. Asymptotes are the
first place a beginner feels that leverage, and it is a feeling worth chasing
through every later chapter.
\end{aestheticbox}

\begin{connectionbox}
\textbf{Asymptotes return as improper integrals.} The vertical and horizontal
asymptotes you sketch here become, a few chapters from now, the boundaries of
\emph{improper integrals} --- areas that stretch out to infinity or pile up
against a cliff. Whether such an infinite region has finite area is decided by
exactly the growth-rate comparisons of this chapter: $\int_1^\infty x^{-2}\,\dd
x$ converges because $1/x^2$ dies fast, while $\int_1^\infty x^{-1}\,\dd x$
diverges because $1/x$ dies too slowly. End behavior is not a side topic; it is
the gatekeeper of which infinite sums and areas make sense.
\end{connectionbox}

\section{Exercises}
\label{ch09:exercises}

\begin{exercise}[Warm-up limits]
Evaluate each limit, naming the dominant terms you used.
\[
\text{(a)}\ \lim_{x\to\infty}\frac{5x^3-2x}{2x^3+7},
\quad
\text{(b)}\ \lim_{x\to\infty}\frac{x+9}{x^2-4},
\quad
\text{(c)}\ \lim_{x\to-\infty}\frac{6x^4+1}{1-3x^2}.
\]
\end{exercise}

\begin{exercise}[Two asymptotes from one radical]
Find both horizontal asymptotes of $f(x)=\dfrac{\sqrt{9x^2+2}}{4x-1}$, treating
$x\to+\infty$ and $x\to-\infty$ separately. Explain, using $\sqrt{x^2}=\abs{x}$,
why the two answers differ in sign.
\end{exercise}

\begin{exercise}[Conjugate trick]
Evaluate $\displaystyle\lim_{x\to\infty}\bigl(\sqrt{x^2+5x}-\sqrt{x^2+x}\bigr)$
by multiplying by the conjugate. (Answer: $2$.) Then explain in one sentence why
the naive answer ``$\infty-\infty=0$'' is wrong.
\end{exercise}

\begin{exercise}[Hole versus asymptote]
For each function, state every vertical asymptote and every hole, with
coordinates.
\[
\text{(a)}\ \frac{x-2}{x^2-4},
\qquad
\text{(b)}\ \frac{x^2+x-6}{x^2-9},
\qquad
\text{(c)}\ \frac{x^2+1}{x^2-2x}.
\]
\end{exercise}

\begin{exercise}[Slant asymptote by division]
Find the slant asymptote of $f(x)=\dfrac{2x^2+3x-1}{x+2}$ and state on which side
of the line the graph lies for large positive $x$ and for large negative $x$.
(Hint: the division leaves $f(x)=2x-1+\frac{1}{x+2}$.)
\end{exercise}

\begin{exercise}[Full sketch]
Carry out the asymptote checklist for $g(x)=\dfrac{x^2-x-2}{x-3}$: holes,
vertical asymptotes, slant asymptote, and all intercepts. Then sketch it,
drawing the asymptotes as dashed lines first.
\end{exercise}

\begin{exercise}[Growth-rate sorting]
Order the following by eventual growth, slowest to fastest, and justify the two
closest calls with a limit:
\[
   x^{50},\qquad (\ln x)^{1000},\qquad 1.01^{x},\qquad \sqrt{x},\qquad x\ln x.
\]
\end{exercise}

\begin{exercise}[A crossing curve]
Give an explicit function with horizontal asymptote $y=2$ that \emph{crosses}
the line $y=2$ infinitely many times as $x\to\infty$. Prove your function does
both. (Hint: a decaying oscillation added to $2$.)
\end{exercise}

\begin{exercise}[Prove a limit from the threshold definition]
Using Definition~\ref{ch09:def-limit-inf} directly, prove that
$\displaystyle\lim_{x\to\infty}\frac{2x-1}{x+4}=2$. That is, given $\eps>0$,
produce an explicit $M$ (in terms of $\eps$) so that $x>M$ forces
$\bigl|\frac{2x-1}{x+4}-2\bigr|<\eps$. (Hint: the difference simplifies to
$\frac{-9}{x+4}$.)
\end{exercise}

\begin{exercise}[Prove log loses to every power]
Adapt the integral bound $\ln x\le 2\sqrt x$ from the proof of
Theorem~\ref{ch09:thm-exp-beats-power} to show that for every fixed $p>0$,
$\displaystyle\lim_{x\to\infty}\frac{\ln x}{x^{p}}=0$. (Hint: it suffices to
beat $x^p$ with $\sqrt x$ when $p\ge \tfrac12$; for small $p$, bound $\ln x$ by
$\tfrac{1}{q}x^{q}$ for a suitable $q<p$ using $\int t^{q-1}\,\dd t$.)
\end{exercise}

\begin{exercise}[Parabolic asymptote]
The function $f(x)=\dfrac{x^3+1}{x}$ has a top exceeding the bottom by two
degrees, so no straight-line asymptote. By dividing, show that $f(x)$ becomes
indistinguishable from a \emph{parabola} for large $\abs{x}$, and state that
parabola. Confirm $f(x)-(\text{parabola})\to 0$.
\end{exercise}

\begin{exercise}[Conceptual]
True or false, with a one-line reason each. (a) A function can have at most two
horizontal asymptotes. (b) A function can have infinitely many vertical
asymptotes. (c) If $\lim_{x\to\infty}f(x)=L$ then $f$ never equals $L$. (d) A
rational function whose denominator has higher degree than its numerator always
has the horizontal asymptote $y=0$.
\end{exercise}
